% Pursuing Stacks (cotes 134-2 to 134-8): archive cote 134-7, site page 108 (PDF page 109), typescript page 550.
% typescript. Draft compared: ps6.tex:5438-5509.
% AI-transcribed by Claude Opus 5.5 (claude-opus-5-5, adaptive thinking, effort high), Message Batches API, batch ps4 (msgbatch_012Aoc1NuskATPsyBeKSYK7J), 2026-11-12; not yet checked by eye.
\scanpage{134-7/108}
\mspage{550}

the three functors $L\overset{\hand{\sim}}{H}.$ , $\overset{\hand{\sim}}{K}$ , $\overset{\hand{\sim}}{S}$ . Indeed, as we just recalled, the
first of the three functors is well defined and has an evident descripti-%
on for all semisimplicial affine schemes over k. As for $\overset{\hand{\sim}}{K}$ and $\overset{\hand{\sim}}{S}$ ,
they are obtained, in terms of \inserted{a variable} co-semisimplicial k-module $L^*$ (represen-%
ting the semisimplicial vector bundle \xout{V(L} $L_* = V(L^*)$) by \xout{taking} \inserted[typed]{applying}
componentwise the functor $\mathrm{Sym}_k$ and the first-order truncation
$\mathrm{Sym}_k(\unsure{;})(1)$ respectively\xout{\unsure{,}}- one may hardly imagine something simpler !
This brings to my attention that in terms of this description, we get
a canonical functorial map
\[
(13)\qquad \overset{\hand{\sim}}{S}(L_*)\hand{\hookrightarrow}\overset{\hand{\sim}}{K}(L_*\hand{)}
\]
when working with the semisimplicial models, and hence presumably a cor-%
responding map for the \inserted{functors between the} relevant derived categories $\mathrm{Hotab}_o(k)$ and $\mathrm{Hot}_o(k)$.
Working either in the model or in the derived categories, this map as
a matter of fact may be deduced from the basic formulaire of section
124, where it had by then escaped my attention. Indeed, \unsure{as} $\overset{\hand{\sim}}{S}$ is a
left adjoint of $L\over{H}{\Pi}.$ , and $\overset{\hand{\sim}}{K}$ a right adjoint of $L\overset{\hand{\sim}}{H}.$, to give such
a map (13) is equivalent with giving either one of two maps
\[
(14)\qquad \left\{
\begin{array}{ll}
a) & L_* \longrightarrow L\over{H}{\Pi}.(\overset{\hand{\sim}}{K}(L_*))\\
b) & L\overset{\hand{\sim}}{H}.(\overset{\hand{\sim}}{S}(L_*)) \longrightarrow L_* \qquad ,
\end{array}\right.
\]
and the formulaire provides for two such maps, namely the ``inversion
isomorphisms''((9), \inserted{p.525).} Thus there is an extra property which was forgotten
in the formulaire, namely that the two maps (13) associated to the two
inversion isomorphisms should be the same. A nicer way then to state
the formulaire is to consider the map (13) above as a basic data, and
say that the two maps in (14) deduced from it by the adjunction property
should be isomorphisms. The situation is reminiscent of the two ways
by which we could obtain the Hurewics map (p.525, C)) - presumably ,
the basic data for the formulaire of section 124 should be the functors
$\overset{\hand{\sim}}{S}$ and $\overset{\hand{\sim}}{K}$ and the map \xout{\illegible{1}} (13) between them, with the property that the
relevant adjoint functors $L\hand{\Pi}.$ and $L\overset{\hand{\sim}}{H}.$ exist, and that the correspondin\unsure{g}
maps in (14) should be isomorphisms, which then will allow to define
a \uline{unique} Hurewicz map $L\hand{\Pi}. \rightarrow L\overset{\hand{\sim}}{H}.$ .

As long as we are sticking to the purely formal aspect, and even
when working in the larger context of semisimplicial affine schemes over
k satisfying merely $X_{\unsure{0}}=k$ , or more generally still, dropping the last
restriction and taking ``k-pointed'' semisimplicial affine schemes instead,
the \inserted{whole} ``four functors formalism'' \inserted{(including even $L\hand{\Pi}.$)} as contemplated in section 124\xout{,} \hand{(}and with
\margin{108}
